\documentclass{article}

\usepackage{amsmath, amssymb, amsthm}
\usepackage[a4paper, margin=1in]{geometry}

\title{Intro to Probabilistic Primality Tests}
\author{sid}
\date{\today}

\newtheorem{theorem}{Theorem}
\newtheorem{lemma}{Lemma}
\newtheorem{example}{Example}

\newcommand{\Z}{\mathbb{Z}}
\newcommand{\bld}[1]{\textbf{#1}}
\newcommand{\itl}[1]{\textit{#1}}
\newcommand{\N}{\mathbb{N}}
\newcommand{\Q}{\mathbb{Q}}
\newcommand{\R}{\mathbb{R}}
\newcommand{\eps}{\varepsilon}
\newcommand{\poly}{\mathrm{poly}}
\newcommand{\polylog}{\mathrm{polylog}}
\newcommand{\bigO}{\mathcal{O}}

\begin{document}

\maketitle

Probabilistic tests are more rigorous than heuristics in that they provide provable bounds on the probability of being fooled by a composite number.

\vspace*{0.8em}

The basic structure of randomized primality tests is as follows:

\begin{enumerate}
  \item Randomly pick a number $a$.
  \item Check equality (corresponding to the chosen test) involving $a$ and the given number $n$. If the equality fails to hold true, then $n$ is a composite number and $a$ is a \itl{witness} for the compositeness, and the test stops.
  \item Get back to the step one until the required accuracy is reached.
\end{enumerate}

After one or more iterations, if $n$ is not found to be a composite number, then it can be declared \bld{probably prime}.
\end{document}